\section*{Capítulo \ref{ch:b3:total-synthesis} --- Síntese total: estratégia e rotas clássicas}

\begin{solution}{exo:b3:total-synthesis:1}
A 90\,\%: $0.9^{10} = 35$\,\% e $0.9^{20} = 12$\,\%. A 80\,\%: $0.8^{10} = 11$\,\% e $0.8^{20} = 1.2$\,\%.
\end{solution}

\begin{solution}{exo:b3:total-synthesis:2}
$5 + 3 + 1 + 4 = 13$ etapas no total; LLS $5 + 1 + 4 = 10$.
\end{solution}

\begin{solution}{exo:b3:total-synthesis:3}
4-(Dimetilamino)butan-2-ona, $\ce{CH3COCH2CH2N(CH3)2}$: o enol da acetona se adiciona ao íon
imínio $\ce{CH2=N+(CH3)2}$.
\end{solution}

\begin{solution}{exo:b3:total-synthesis:4}
Desconecte as duas ligações do anel alílicas à C=C: buta-1,3-dieno e but-3-en-2-ona
(metil vinil cetona), que reagem por aquecimento.
\end{solution}

\begin{solution}{exo:b3:total-synthesis:5}
Linear: $0.9^{21} = 11$\,\%. Convergente: $0.9^{11} = 31$\,\%, quase três vezes mais.
\end{solution}

\begin{solution}{exo:b3:total-synthesis:6}
A: $(7 + 1)/12 = 67$\,\%. B: $7/9 = 78$\,\%.
\end{solution}

\begin{solution}{exo:b3:total-synthesis:7}
Sem: $0.9^8 = 43$\,\%. Com quatro etapas a mais a 95\,\%: $43\,\% \times 0.95^4 = 35$\,\%.
\end{solution}

\begin{solution}{exo:b3:total-synthesis:8}
Primário: um silil éter volumoso (\textit{terc}-butildifenilsilila), removido com fluoreto.
Secundário: um éter 4-metoxibenzílico, removido por oxidação com DDQ. O álcool
terciário, impedido, pode muitas vezes ser deixado livre, ou protegido como um pequeno silil éter
clivado por ácido brando. Cada condição deixa os outros grupos intactos.
\end{solution}

\begin{solution}{exo:b3:total-synthesis:9}
As sínteses totais exigiam dezenas de etapas, com \omterm{def:b3:total-synthesis:architecture}{rendimentos globais} muito abaixo de 1\,\%; um
composto com o sistema de anéis completo e a maioria dos estereocentros, extraído das
folhas do teixo europeu, uma fonte renovável, é convertido em poucas etapas.
\end{solution}

\begin{solution}{exo:b3:total-synthesis:10}
Adição de Michael: o enolato da diona (C2, entre as carbonilas) se adiciona ao
carbono $\beta$ da metil vinil cetona, dando uma tricetona. Aldólica: o enolato da metilcetona
ataca uma carbonila do anel de forma intramolecular, e a desidratação dá a
ciclo-hexenona: a cetona de Wieland--Miescher, uma enodiona bicíclica com um grupo metila
angular.
\end{solution}

\begin{solution}{exo:b3:total-synthesis:11}
$y^{m+1}/y^{2m+1} = y^{-m}$: a 90\,\%, 1.7 para $m = 5$ e 4.9 para $m = 15$. Quanto mais longa a rota, mais
a convergência compensa.
\end{solution}

\begin{solution}{exo:b3:total-synthesis:12}
Cada anel se fecha com o cátion e a ligação dupla atacante em um arranjo do tipo
cadeira; a adição a cada ligação dupla é anti, e os substituintes ocupam
posições pseudoequatoriais, o que torna \textit{trans} todas as fusões de anéis. Uma ligação dupla \textit{E}
dispõe suas duas continuações de cadeia de modo que resulte a fusão \textit{trans}; uma \textit{Z}
daria a fusão \textit{cis}: a geometria do polieno codifica a estereoquímica do
produto.
\end{solution}

\begin{solution}{pb:b3:total-synthesis:1}
\textbf{1.} $\ce{C4H6O2}$, $\ce{CH5N}$, $\ce{C5H6O5}$, $\ce{C8H13NO}$.

\textbf{2.} $\ce{C4H6O2 + CH5N + C5H6O5 -> C8H13NO + 2CO2 + 2H2O}$.

\textbf{3.} 86.0, 31.0 e \qty{146.0}{g/mol}; tropinona \qty{139.0}{g/mol}.

\textbf{4.} $139.0/(86.0 + 31.0 + 146.0) = 0.53$: 53\,\%.

\textbf{5.} Seus grupos $\ce{CH2}$ centrais, cada um entre dois grupos carbonila, se enolizam facilmente em
água; a própria acetona é muito pouco nucleofílica, e os grupos carboxila ativantes
saem depois.

\textbf{6.} Os íons imínio exigem amina livre (ácido demais a protona), e o
enol e a formação do imínio exigem algum ácido: um pH perto da neutralidade os equilibra.

\textbf{7.} O $\ce{CH3NH2}$ se adiciona a um grupo aldeído; a perda de água dá o íon imínio
$\ce{R-CH=N+H-CH3}$.

\textbf{8.} Um carbono enólico do ácido acetonodicarboxílico ataca o carbono do imínio, formando uma
ligação C--C e uma amina secundária.

\textbf{9.} A amina secundária se condensa com o segundo grupo aldeído da mesma
cadeia, dando um íon imínio cíclico (de cinco membros).

\textbf{10.} O enol do outro lado da cetona o ataca, fechando o
anel de seis membros e o esqueleto bicíclico.

\textbf{11.} Cada um é um $\beta$-cetoácido, que perde $\ce{CO2}$ por um estado de transição cíclico de seis
membros, dando o enol da cetona.

\textbf{12.} Duas ligações C--C e duas ligações C--N.

\textbf{13.} $0.75^{14} = 0.018$: 1.8\,\%.

\textbf{14.} $42/1.78 = 24$.

\textbf{15.} 100\,\%: uma única etapa, toda de construção de ligações.

\textbf{16.} Uma única operação, água como solvente, componentes baratos, nenhum grupo protetor e
quase nenhum resíduo além de dióxido de carbono e água.

\textbf{17.} A planta constrói o esqueleto tropano a partir de um íon imínio cíclico derivado de um
aminoácido e de uma unidade derivada do acetato, pelo mesmo tipo de química de
Mannich.

\textbf{18.} As cascatas (e as \omterm{def:b3:pericyclic:families}{cicloadições} formadoras de anéis, como a reação de
Diels--Alder).

\textbf{19.} Não: um \omterm{def:b3:point-groups:mirror}{plano especular} passa por N, por seu grupo metila, por C3 e pelo oxigênio
carbonílico.

\textbf{20.} Cada um tem quatro grupos diferentes, mas eles são imagens especulares um do outro pelo
plano de simetria da própria molécula: um composto do tipo meso.

\textbf{21.} Diastereoisômeros, que diferem pela orientação do OH em C3 em relação à
ponte de nitrogênio.

\textbf{22.} Não: o \omterm{def:b3:point-groups:mirror}{plano especular} se conserva nos dois.

\textbf{23.} As duas faces do grupo carbonila não são equivalentes: uma está voltada para a
ponte de nitrogênio, a outra para a ponte de dois carbonos, que impedem de modo diferente.

\textbf{24.} A rota em um só pote dá cerca de 24 vezes o \omterm{def:b3:total-synthesis:architecture}{rendimento global} da rota
linear de 14 etapas.
linear route.
\end{solution}
